\documentclass[a4paper,11pt]{article}
\usepackage[utf8]{inputenc}
\usepackage[T1]{fontenc}
\usepackage[french]{babel}
\usepackage{lmodern}
\usepackage{amsmath,amssymb,amsthm}
\usepackage{mathrsfs}
\usepackage{enumitem}
\usepackage{multicol}

\setlist[enumerate]{label=\textbf{\textcolor{Couleur8}{\arabic*.}}, left=1em}

% Si vous voulez aussi modifier les listes enumerate imbriquées :
\setlist[enumerate,2]{label=\textcolor{Couleur8}{\alph*)}}
\setlist[enumerate,3]{label=\textcolor{Couleur8}{\roman*)}}
\setlist[enumerate,4]{label=\textcolor{Couleur8}{\Alph*)}}
\usepackage{graphicx}
\usepackage{xcolor}
\usepackage{tcolorbox}
\usepackage{geometry}
\usepackage{fancyhdr}
\usepackage{titlesec}
\usepackage{cancel}
\usepackage{ifthen}
\usepackage{tikz}
\usetikzlibrary{arrows.meta}
\usepackage{tkz-tab}
\usepackage{eurosym}

% OPTION PROF/ÉLÈVE
% Décommentez UNE des deux lignes suivantes :
\newboolean{prof}\setboolean{prof}{true}   % Version professeur (avec corrections des devoirs)
\newboolean{prof}\setboolean{prof}{true}  % Version élève (sans corrections des devoirs)

\definecolor{Couleur1}{HTML}{4A5D7A} %Th
\definecolor{Couleur2}{HTML}{A5B5D6} %Prop
\definecolor{Couleur3}{HTML}{A8C68A} %Méthode
\definecolor{Couleur6}{HTML}{8A9CC1} %Def
\definecolor{Couleur7}{HTML}{E6B459} %Rq Voc
\definecolor{Couleur8}{HTML}{D36740} %Titres
\definecolor{correction}{HTML}{D36740} % Couleur pour les corrections
\definecolor{coopmathscorr}{HTML}{F15929} % Couleur corrections coopmaths

% Configuration des marges
\geometry{
	top=2cm,
	bottom=2cm,
	inner=1.5cm,
	outer=1.5cm,
}

% Police
\renewcommand{\familydefault}{\sfdefault}

% Compteurs
\newcounter{seancenum}
\newcounter{seancenumD}
\setcounter{seancenum}{1}
\setcounter{seancenumD}{1}

% Configuration des en-têtes et pieds de page
\pagestyle{fancy}
\fancyhf{}
\ifthenelse{\boolean{prof}}{
    \fancyhead[L]{\textcolor{Couleur8}{\textbf{Corrections Automatismes - Classe de seconde - Version Professeur}}}
}{
    \fancyhead[L]{\textcolor{Couleur8}{\textbf{Corrections Devoirs - Séance 23 - Classe de seconde}}}
}
\fancyhead[R]{\textcolor{Couleur8}{\textbf{2025-2026}}}
\fancyfoot[L]{\textcolor{Couleur8}{Mme Bertail et Mme Pinto}}
\fancyfoot[R]{\textcolor{Couleur8}{\thepage}}
\renewcommand{\headrulewidth}{1pt}
\renewcommand{\headrule}{\hbox to\headwidth{\color{Couleur8}\leaders\hrule height \headrulewidth\hfill}}

% Environnements personnalisés
\newtcolorbox{seance}[1]{
    colback=Couleur6!10,
    colframe=Couleur1!65,
    fonttitle=\bfseries\large,
    title=Correction Séance \theseancenum #1,
    left=5mm,
    right=5mm,
    top=3mm,
    bottom=3mm,
    arc=0mm,
    after=\stepcounter{seancenum}
}

\newtcolorbox{seanceD}[1]{
    colback=Couleur6!5,
    colframe=Couleur6!45,
    fonttitle=\bfseries\large,
    title=Correction Devoirs - Séance \theseancenumD #1,
    left=5mm,
    right=5mm,
    top=3mm,
    bottom=3mm,
    arc=0mm,
    after=\stepcounter{seancenumD}
}

\begin{document}

\setcounter{seancenumD}{23}
\newpage
\begin{seanceD}{} %23

\begin{enumerate}
\item $4 \times 1{,}5={\color[HTML]{f15929}\boldsymbol{6}}$
	\item $45-50+8={\color[HTML]{f15929}\boldsymbol{3}}$
	\item $
      (x+1)(x+4)=x^2+4x+x+4
      ={\color[HTML]{f15929}\boldsymbol{x^2+5x+4}}$\\Le terme en $x^2$ vient de $x\times 1x=x^2$.\\Le terme en $x$ vient de la somme de $x \times 4$ et de $1 \times 1x$.\\Le terme constant vient de $1\times 4= 4$.
	\item $
      2+\dfrac{5}{3} = \dfrac{2 \times 3}{3} + \dfrac{5}{3} = \dfrac{6}{3} + \dfrac{5}{3}  ={\color[HTML]{f15929}\boldsymbol{\dfrac{11}{3}}}$
	\item $20\,\%$ de $40 = {\color[HTML]{f15929}\boldsymbol{8}}$\\ Prendre $20\,\%$  de $40$ revient à prendre $2\times 10\,\%$  de $40$.\\
      Comme $10\,\%$  de $40$ vaut $4$ (pour prendre $10\,\%$  d'une quantité, on la divise par $10$), alors
      $20\,\%$ de $40=2\times 4=8$.
     
	\item On décompose le calcul pour le rendre plus simple mentalement :\\  $ 0{,}4\times0{,}5 =4\times 0,1\times 5\times 0,1= 20\times 0,01
= {\color[HTML]{f15929}\boldsymbol{0{,}2}}$
	\item Comme $1{,}11=1+0{,}11$, multiplier par $1{,}11$ revient à augmenter de ${\color[HTML]{f15929}\boldsymbol{11}}\,\%$. 
	\item La somme des $3$ nombres est : $4+12+11 =27$.\\
            La moyenne est donc $\dfrac{27}{3}={\color[HTML]{f15929}\boldsymbol{9}}$.
	\item $\sqrt{49}={\color[HTML]{f15929}\boldsymbol{7}}$
	\item  L'algorithme retourne $2+(-4)={\color[HTML]{f15929}\boldsymbol{-2}}$. 
	\item Le nombre qui multiplié par $6$ donne $13$ est ${\color[HTML]{f15929}\boldsymbol{\dfrac{13}{6}}}$.
      
	\item L'abscisse du point $A$ est $\dfrac{6}{4}={\color[HTML]{f15929}\boldsymbol{1{,}5}}$.
	\item L'aire est donnée par le produit des deux plus petits côtés divisé par $2$, soit : $\dfrac{3\times 4}{2}={\color[HTML]{f15929}\boldsymbol{6}}\text{ cm}^2$.
	\item  $10\,\%$ de  $1\,000$\,\euro{}~est égal à $100$\,\euro{}. \\
        Après la hausse de $10\,\%$, le prix est de $1\,000$\,\euro{}~$+$ $100$\,\euro{}~$=1\,100$\,\euro{}.\\
        $10\,\%$ de  $1\,100$\,\euro{}~est égal à $110$\,\euro{}. \\
        Après la baisse de $10\,\%$, le prix est de  $1\,100$\,\euro{}~$-$ $110$\,\euro{}~$=990$\,\euro{}.\\
        Le nouveau prix est ${\color[HTML]{f15929}\boldsymbol{990}}$\,\euro{}. \\
        \textbf{Autre méthode :\\ }
        Augmenter de $10\,\%$ revient à multiplier par $1{,}1$.\\
Diminuer de $10\,\%$ revient à multiplier par $0{,}9$.\\
Donc le prix est multiplié par $1{,}1\times 0{,}9$, c'est-à-dire par $0{,}99$.\\
Le prix final est donc inférieur au prix initial (car on multiplie par un nombre inférieur à $1$).
        
	\item $
      10^{-2}+10^{3}=0{,}01 +1\,000
      ={\color[HTML]{f15929}\boldsymbol{1\,000{,}01}}$

\end{enumerate}


\end{seanceD}

\end{document}